\section{Materials}
In this section we provide a brief overview of the materials used in this project. 

\subsection{Dataset}
The data used in this project originates from the Wikispeedia project \cite{semdistances,humanwayfinding}. The dataset and more information can be found at \href{https://dlab.epfl.ch/wikispeedia/play/}{https://dlab.epfl.ch/wikispeedia/play/}.

The following items are relevant for us in this paper:
\begin{itemize}
    \item articles: a list of all articles
    \item links: a list of all articles with a link between them
    \item paths (un)finished: a collection of paths traversing the graph 
\end{itemize}

The graph in this case is defined over the articles as nodes and links as edges. 

\subsection{Transformer}
To create our embeddings we use a BERT based sentence-encoder. Specifically we use a version of the MiniLLM model\cite{gu2024minillmknowledgedistillationlarge} which can be found on Hugging Face (a repository for machine learning resources) with the following identifier: \textit{sentence-transformers/all-MiniLM-L6-v2}. It maps sentences and paragraphs to a 384 dimensional dense vector space.